%% Īpašie uzstādījumi
%% (kas ir "īpašs" gan ir izvēlēts patvaļīgi) :)

%% Pakotnes 'listings' koda fragmentu konfigurācija
\lstset{%
	% Ar rindiņu numuriem
	numbers=left, numberstyle=\tiny, numbersep=5pt,
	% <tab> = 2 simboli
	tabsize=2,%
	% "Ierāmēt" ar horizontālām līnijām
	frame=lines,framerule=0.5pt,%
	%belowcaptionskip=2ex,%
	% Šrifts: 10pt Monospace
	basicstyle=\footnotesize\ttfamily,%
	% Rindiņas ciešāk (jo \onehalfspacing ir aktīvs)
	lineskip=-0.5ex,
	% Sintakses krāsu un stila izvēles
	keywordstyle={\bfseries\color{blue}},%
	keywordstyle=[2]{\color{Brown}},%
	keywordstyle=[3]{\color{OliveGreen}},%
	keywordstyle=[4]{\color{TealBlue}},%
	keywordstyle=[5]{\color{Sepia}},%
	emphstyle={\color{Bittersweet}},%
	emphstyle={[2]\color{PineGreen}},%
	emphstyle={[3]\color{Green}},%
	commentstyle=\slshape\color{gray},
	stringstyle=\color{orange},
	% Ar šo var izcelt globāli izmantotos identifikatorus
	emph={[3]{cpu_lib,bus_t,t_alu,t_reg,t_shift}},% My global emphs
}

%% Pirmkoda valodu definīcijas
%% Ja 'listings' neatbalsta tavējo, definē (vai papildini) jaunas!
\lstdefinelanguage[qucs]{VHDL}[]{VHDL}%
{morekeywords=[1]{rising_edge},% MOAR core keywords
morekeywords=[3]{ieee,std_logic_1164,std_logic_arith,std_logic_unsigned,%
	work,std,textio,% usual libraries
	std_logic,std_logic_vector,time,integer,real,line,string,% data types
	ns},%real type units
morekeywords=[4]{reverse_range,length,event,last_value,high,low},% attributes
morekeywords=[5]{write,writeline,conv_integer},%funkcijas
morestring=[b]',
}

\input{jsi_assembler_listing_def.tex}

%% Identifikatori tekstā? Jā, var!
%% Lietot prātīgi! "With great power, comes great responsibility!"
\lstMakeShortInline[language={[qucs]VHDL},basicstyle=\normalsize\ttfamily]|

%% English determined number suffixes
%\renewcommand{\th}{\textsuperscript{th} }
\newcommand{\st}{\textsuperscript{st} }
\newcommand{\nd}{\textsuperscript{nd} }
\newcommand{\rd}{\textsuperscript{rd} }
\newcommand{\nth}{\textsuperscript{th} }

%% Anotācijas virsraksta stils
\newcommand{\abstitlestyle}[1]{%
	\noindent \begin{center}
		\textbf{\Large #1}
	\end{center}}
